
\documentclass[12pt,a4paper,pl,fleqn]{article}


% \usepackage[english,polish]{babel}                   % ostatni jezyk przewaza w opisach rozdzialow i literatury
\usepackage[utf8]{inputenc}
\usepackage{graphicx}
\usepackage[toc]{appendix}
% \usepackage[T1]{fontenc}



% \documentclass[a4paper,12pt]{article}
% \usepackage[utf8]{inputenc}
\usepackage[T1]{fontenc}
\usepackage[polish]{babel}
\usepackage{polski}
\usepackage{lmodern}


\setlength{\textheight}{210mm}
\setlength{\textwidth}{145mm}
\setlength{\parindent}{20.0mm}

\renewcommand{\thesection}{\Roman{section}}

\title{Równanie oscylatora harmonicznego prostego z rozwiązaniem}

\author{Paweł Jan Piskorz}



\date{}

%++++++++++++++++++++++++++++++++++++++++++++++++++++++++++
\begin{document}


\maketitle

Rozważmy zastosowanie drugiej zasady dynamiki Newtona kiedy siła działająca
na punkt materialny jest proporcjonalna do jego wychylenia z położenia
równowagi i skierowana do położenia równowagi. Jako przykład rozważmy jednowymiarowy ruch
wzdłuż osi $Ox$\@. Wówczas mamy następujące równanie ruchu Newtona
% We will consider the application of the Newton second law when the force acting on a material point
% is proportional to the displacement from the equilibrium point and directed
% to the equilibrium point.    % ~\cite{kn:Guminski_Book}.
% As an example let us consider a one dimensional motion along the $x$ axis. The Newton equation then is
\begin{equation}
m \ddot{x} = - k x
\end{equation}
gdzie $x$ jest wychyleniem, $m$ jest masą, $k$ jest stałą proporcjonalności zwaną stałą siłową natomiast
% where $x$ is displacement, $m$ is mass, $k$ is proportionality constant, and
\begin{equation}
\ddot{x} = \frac{d^2 x}{d t^2}
\end{equation}
jest drugą pochodną po czasie $t$ wychylenia $x$ oscylatora czyli jest przyspieszeniem.

Całkiem podobne równanie otrzymujemy rozważając wahadło o masie $m$ zawieszonej na długiej nierozciągliwej nici o długości $l$
w przybliżeniu wychyleń o niewielki kąt co odpowiada wychyleniom z położenia równowagi o niewielką odległość $x$
% is the second time $t$ derivative of the displacement, i.e.~the accelleration.
% Quite similar equation we get when treating the long pendulum in small displacement angle approximation
\begin{equation}
m \ddot{x} = - m g (x / l)
\end{equation}
% what gives
co daje
\begin{equation}
l \ddot{x} = - g x
\end{equation}
% where $l$ is the length of the pendulum, $g$ is the acceleration due to gravity.
gdzie $g$ jest przyspieszeniem ziemskim.
Wahadło o powyższym równaniu ruchu nazywamy {\em wahadłem matematycznym.}

% Our goal will be to solve the equation
Naszym celem będzie rozwiązanie równania różniczkowego rzędu drugiego o stałych współczynnikach zwanego
{\em równaniem oscylatora harmonicznego prostego\/}
\begin{equation}
\label{eq:harmonic_oscillator_equation}
m \frac{d^2 x}{d t^2} + k x = 0
\end{equation}
bez zgadywania rozwiązania.
% without making any assumptions about the solution function $x$\@.    % as suggested in~\cite{kn:Romanowski_Wrona_Book}.
% We will make assumptions only about the initial conditions.
% The velocity $v$ of the motion is
Uczynimy jedynie założenia dotyczące warunków początkowych.

Prędkość $v$ ruchu oscylatora możemy zapisać jako
\begin{equation}
\dot{x} = v = \frac{d x}{d t}
\end{equation}

% We start with noticing that
Rozpoczynamy zauważając, że
\begin{equation}
\frac{d^2 x}{d t^2} = \frac{d}{d t} \left( \frac{d x}{d t} \right)  =  \frac{d v}{d t}  =   
\frac{d v}{d x} \frac{d x}{d t}
=    \frac{d v}{d x} v                             
\end{equation}
\begin{equation}
     \frac{d (v^2)}{d x}     =   2 v \frac{d v}{d x}  
\end{equation}
\begin{equation}
      v \frac{d v}{d x}   =  \frac{1}{2} \frac{d (v^2)}{d x} =  \frac{d^2 x}{d t^2}      
\end{equation}
% The original equation (\ref{eq:harmonic_oscillator_equation}) becomes
Wyjściowe równanie (\ref{eq:harmonic_oscillator_equation}) przekształca się w równanie
\begin{equation}
     m \frac{1}{2} \frac{d (v^2)}{d x} +   k x = 0
\end{equation}
\begin{equation}
     d (v^2)   +   \frac{k}{m} 2 x dx = 0
\end{equation}
% After integration
Po scałkowaniu powyższego równania dostajemy
\begin{equation}
     v^2  +  \frac{k}{m} x^2  +  C_{1}   =    0
\end{equation}
% where $C_{1}$ is a constant. For maximum displacement $x = x_{0}$ the velocity $v = 0$ and
gdzie $C_{1}$ jest stałą. Dla maksymalnego wychylenia $x = x_{0}$ prędkość $v = 0$ i
\begin{equation}
      \frac{k}{m} x_{0}^2  +  C_{1}   =    0
\end{equation}
\begin{equation}
      C_{1} = - \frac{k}{m} x_{0}^2
\end{equation}
% what gives
co daje
\begin{equation}
      v^2  = \frac{k}{m} \left( x_{0}^2 - x^2 \right)
\end{equation}
\begin{equation}
      \frac{d x}{d t}  =   \sqrt{\frac{k}{m}} x_{0} \sqrt{ 1 - \left( \frac{x}{x_{0}} \right)^2 }
\end{equation}
% Substituting $u = x / x_{0}$ and $dx = x_{0} du$
Podstawiając $u = x / x_{0}$ and $dx = x_{0} du$
\begin{equation}
\label{eq:integration_for_arcsin}
       \sqrt{\frac{k}{m}} \,  d t  =  \frac{d x}{ x_{0} \sqrt{ 1 - \left( x/x_{0} \right)^2 } }
       = \frac{x_{0} du}{ x_{0} \sqrt{ 1 - \left( x/x_{0} \right)^2 } }
       =  \frac{d u}{ \sqrt{ 1 - u^{2}} }
\end{equation}
% Integrating the equation (\ref{eq:integration_for_arcsin}) we receive
Całkując równanie (\ref{eq:integration_for_arcsin}) otrzymujemy
\begin{equation}
    \arcsin u =  \arcsin(x/x_{0}) =  \sqrt{k/m} \,   t + C_{2}
\end{equation}
% and from that
a z tego
\begin{equation}
    x = x_{0} \sin(\sqrt{k/m} \,   t + C_{2})
\end{equation}
% We may assume the that the constant $C_{2} = 0$ if for $t = 0$ the displacement $x = 0$, and our solution
% of the equation (\ref{eq:harmonic_oscillator_equation}) is
Możemy założyć, że stała $C_{2} = 0$ jeżeli dla $t = 0$ wychylenie $x = 0$, i naszym poszukiwanym rozwiązaniem
równania (\ref{eq:harmonic_oscillator_equation}) jest
\begin{equation}
\label{eq:harmonic_oscillator_equation_solution}
    x = x_{0} \sin(\sqrt{k/m} \, t)
\end{equation}
% Because of the above solution the equation (\ref{eq:harmonic_oscillator_equation}) is called
% the equation of a harmonic oscillator.

% The solution of equation (\ref{eq:harmonic_oscillator_equation})
% is a sinusoidal function. Sinusoidal function also describes the vibration of a string e.g.~of a violin
% or a harp string. Such vibration corresponds to a harmonic tone and therefrom comes the name harmonic oscillation.
% Because of its solution (\ref{eq:harmonic_oscillator_equation_solution})
% the equation (\ref{eq:harmonic_oscillator_equation}) is called
% the equation of a harmonic oscillator.


Rozwiązaniem równania (\ref{eq:harmonic_oscillator_equation}) oscylatora harmonicznego
jest funkcja sinusoidalna (\ref{eq:harmonic_oscillator_equation_solution})\@. 
Funkcja sinusoidalna opisuje również drgania strun, na przykład skrzypcowych, gitary czy harfy.
Drgania takie odpowiadają tonom harmonicznym i stąd bierze się nazwa {\em drganie harmoniczne oscylatora.\/}
Zasadniczo każdy punkt drgającej struny oscyluje jak oscylator harmoniczny.



% Essentially, every point on a vibrating string vibrates like a harmonic oscillator.


% From the periodicity of the sinusoidal function
% we see then that the motion of the harmonic oscillator is periodic. Let us denote the period of the
% motion of the harmonic oscillator by $T$. Then


Z okresowości funkcji sinusoidalnej widzimy zatem, że ruch oscylatora harmonicznego jest okresowy.
Oznaczmy okres ruchu oscylatora harmonicznego przez $T$\@. Wtedy
\begin{eqnarray}
    x(t) = x_{0} \sin(\sqrt{k/m} \, t)
    = x_{0} \sin(\sqrt{k/m} \, t + 2\pi)  \\
    = x_{0} \sin(\sqrt{k/m} \, ( t + T ))
    = x_{0} \sin(\sqrt{k/m} \, t + \sqrt{k/m} \, T)  \nonumber
\end{eqnarray}
% and from that
i z powyższego
\begin{equation}
    T \sqrt{k/m} = 2 \pi
\end{equation}
% with the period $T$
mamy ruch z okresem $T$
\begin{equation}
    T  = 2 \pi \sqrt{m/k}
\end{equation}
Po czasie $T$ ta sama faza ruchu oscylatora powtarza się. Przez fazę ruchu oscylatora
rozumiemy pewien określony stopień zaawansowania jego ruchu. Mamy na przykład fazę pełnego wychylenia w prawo,
 fazę przejścia przez położenie równowagi z prawej strony w lewo, itp.

% We notice that for the pendulum we have the period
Zauważamy, że dla wahadła matematycznego mamy okres
\begin{equation}
    T  = 2 \pi \sqrt{l/g}
\end{equation}

% Frequency $\nu$ of the periodic motion denotes the number of periods per unit time
Częstotliwość $\nu$ ruchu okresowego oznacza ilość okresów ruchu przypadających na jednostkę czasu
\begin{equation}
\label{eq:frequency}
    \nu = 1/T
\end{equation}
Mogliśmy zapisać jak powyżej ponieważ na czas $T$ przypada dokładnie jeden pełny okres ruchu oscylatora.

Częstotliwość kołowa $\omega$ ruchu okresowego to z definicji
\begin{equation}
\label{eq:circular_frequency}
    \omega = 2 \pi/T
\end{equation}
Mamy, że
\begin{equation}
\label{eq:circular_frequency_explained}
     \omega = \frac{2\pi}{T} = \sqrt{k/m}
\end{equation}


Często w podręcznikach równanie ruchu oscylatora harmonicznego prostego zapisane jest od razu w postaci
\begin{equation}
\ddot{x} = - {\omega}^{2} x
\end{equation}
a rozwiązanie powyższego równania to
\begin{equation}
x(t) = x_{0} \sin(\omega \, t)
\end{equation}
Prędkość oscylatora otrzymujemy z obliczenia pochodnej wychylenia po czasie
\begin{equation}
v(t) = \frac{d x(t)}{d t} = \frac{d}{d t} x_{0} \sin(\omega \, t) = x_{0} \omega \cos(\omega \, t) 
= v_{0} \cos(\omega \, t)
\end{equation}







% For the harmonic oscillator
Dla oscylatora harmonicznego prostego
\begin{equation}
\label{eq:frequency_of_harmonic_ooscillator}
    \nu = \frac{1}{2 \pi} \sqrt{k/m}
\end{equation}
% and for the pendulum in the approximation of harmonic oscillations
a dla wahadła matematycznego
\begin{equation}
\label{eq:frequency_of_pendulum}
    \nu = \frac{1}{2 \pi} \sqrt{g/l}
\end{equation}

% So our simple harmonic oscillator problem is solved. We had to know that


Dodatkowo rozważmy energię ruchu oscylatora harmonicznego. Energia ta jest sumą energii kinetycznej
\begin{equation}
    \frac{1}{2} m v^{2} 
\end{equation}
oraz energii potencjalnej
\begin{equation}
    \frac{1}{2} k x^{2} 
\end{equation}
Energię potencjalną ruchu oscylatora oblicza się z całki pracy
\begin{equation}
    -\int_{0}^{x} \vec{F} \cdot \vec{i} ds =
    -\int_{0}^{x} -ks \, ds =  \frac{1}{2} k x^{2}
\end{equation}
gdzie $\vec{F}$ jest wektorem siły działającej na punkt materialny oscylatora, $s$ jest przesunięciem
natomiast $\vec{i}$ jest wersorem osi $Ox$ układu współrzędnych.

Energia całkowita ruchu oscylatora harmonicznego jest sumą jego e\-ner\-gii kinetycznej oraz potencjalnej
\begin{eqnarray}
    \frac{1}{2} m v^{2} + \frac{1}{2} k x^{2} =                     \\ \nonumber
  = \frac{1}{2} m v_{0}^{2} \cos^{2}(\omega \, t) +
    \frac{1}{2} k x_{0}^{2} \sin^{2}(\omega \, t)
\end{eqnarray}
Ponieważ $v_{0} = x_{0} \omega$ to 
\begin{equation}
m v_{0}^{2} = m x_{0}^{2} \omega^{2} = m x_{0}^{2} \frac{k}{m} = k x_{0}^{2}
\end{equation}
i otrzymujemy, że energia całkowita $E$ oscylatora to
\begin{equation}
E = \frac{1}{2} m \frac{k}{m} x_{0}^{2} \cos^{2}(\omega \, t) + \frac{1}{2} k x_{0}^{2} \sin^{2}(\omega \, t)
=
\frac{1}{2} k x_{0}^{2}
\end{equation}

Zatem nasz problem znalezienia rozwiązania równania ruchu os\-cy\-la\-to\-ra
harmonicznego prostego został rozwiązany. Jedynie musieliśmy znać całkę
\begin{equation}
\label{eq:integral_for_arcsin}
    \int \frac{d u}{ \sqrt{ 1 - u^{2}} } = \arcsin u + C
\end{equation}
% where $C$ is constant. Let us prove the above equation (\ref{eq:integral_for_arcsin}) by asking how to compute
% the derivative of the $\arcsin$ function.
gdzie $C$ oznacza stałą. Udowodnijmy powyższe równanie (\ref{eq:integral_for_arcsin}) zadając pytanie,
w jaki sposób obliczyć pochodną funkcji $\arcsin u$\@.

% Let
Niech
\begin{equation}
                y = \sin x
\end{equation}
\begin{equation}
                \frac{d y}{d x} = \frac{d \sin x}{d x} = \cos x = \sqrt{1 - \sin^2 x}
\end{equation}
% With     
Z 
$y = \sin x$  i   $x = \arcsin y$ stosujemy wzór na obliczenie pochodnej funkcji odwrotnej
% we use the formula to calculate the derivative of the inverse function
\begin{equation}
                \frac{d x}{d y} =  \frac{1}{\sqrt{1 - \sin^2 x}} = \frac{d \arcsin y}{d y}    =   \frac{1}{\sqrt{1 - y^2}}
\end{equation}
% what is equivalent to
co jest równoważne
\begin{equation}
\label{eq:arcsin_derivative}
                \frac{d \arcsin u}{d u} = \frac{1}{\sqrt{1 - u^2}}
\end{equation}
% and the equation (\ref{eq:integral_for_arcsin}) can be proven
% after integrating (\ref{eq:arcsin_derivative})
i równanie (\ref{eq:integral_for_arcsin}) może być dowiedzione
po scałkowaniu równania (\ref{eq:arcsin_derivative})
\begin{equation}
\label{eq:arcsin_derivative_integration}
                \int \frac{1}{\sqrt{1 - u^2}} \, du
                =
                \int d \, \arcsin u
                =
                \arcsin u + C
\end{equation}

%8888888888888888888888888888888888888888888888888888888888888888

% \begin{thebibliography}{99}
% \bibitem{kn:Guminski_Book} Guminski,~K.\
%                           (1964) {\em Elementy chemii teoretycznej (Elements of Theoretical Chemistry)\/} Panstwowe Wydawnictwo Naukowe (State Scientific Publishing), Warszawa
% \bibitem{kn:Romanowski_Wrona_Book} Romanowski,~S., Wrona,~W.\
%                           (1967) {\em Matematyka wyzsza dla studiow technicznych (Higher Mathematics for Technical Studies)\/} Panstwowe Wydawnictwo Naukowe (State Scientific Publishing), Warszawa

% \end{thebibliography}

% \subsection*{}

% Pawel Jan Piskorz (paweljanpiskorz@gmail.com)

\end{document}




